1. even if sometimes we cannot encode the whole decision tree, but it might have very succinct form (say, a simple circuit) which has exactly the same computational result. So even if we can prove tree size large, doesn't means we cannot find a efficient algorithm. ref: DGV08 example 1

2. It seems that in paper Approximation Algorithms for Optimal Decision trees and Adaptive TSP Problems they have a different def with us. Our def of adaptive is more like stochastic probing.